\documentclass[10pt, a4paper, twocolumn]{article}

% --- PAQUETES DE FORMATO Y DISEÑO ---
\usepackage[utf8]{utf8x}
\usepackage[margin=1cm, top=1.2cm, bottom=1.2cm]{geometry}
\usepackage{listings}
\usepackage{xcolor}
\usepackage{amsmath, amssymb}
\usepackage{multicol}
\usepackage{enumitem}
\usepackage{titlesec}
\usepackage{hyperref}

% --- CONFIGURACIÓN DE PÁGINA Y FUENTES ---
\pagestyle{plain}
\setlist{noitemsep, topsep=0pt, parsep=0pt, partopsep=0pt}

% --- CONFIGURACIÓN DE LISTINGS (C++) ---
\definecolor{codegreen}{rgb}{0,0.6,0}
\definecolor{codegray}{rgb}{0.5,0.5,0.5}
\definecolor{codepurple}{rgb}{0.58,0,0.82}
\definecolor{backcolour}{rgb}{0.97,0.97,0.97}

\lstdefinestyle{cppstyle}{
    backgroundcolor=\color{backcolour},   
    commentstyle=\color{codegreen},
    keywordstyle=\color{blue},
    numberstyle=\tiny\color{codegray},
    stringstyle=\color{codepurple},
    basicstyle=\ttfamily\tiny,
    breakatwhitespace=false,         
    breaklines=true,                 
    captionpos=b,                    
    keepspaces=true,                 
    numbers=left,                    
    numbersep=4pt,                  
    showspaces=false,                
    showstringspaces=false,
    showtabs=false,                  
    tabsize=2
}
\lstset{style=cppstyle}

% --- FORMATO DE TÍTULOS ---
\titleformat{\section}{\large\bfseries}{\thesection}{0.5em}{}
\titleformat{\subsection}{\normalfont\bfseries}{\thesubsection}{0.5em}{}
\titleformat{\subsubsection}{\small\bfseries}{\thesubsubsection}{0.5em}{}
\titlespacing*{\section}{0pt}{4pt}{2pt}
\titlespacing*{\subsection}{0pt}{3pt}{1pt}
\titlespacing*{\subsubsection}{0pt}{2pt}{1pt}

\begin{document}

\title{\textbf{ACM-ICPC Cheat Sheet (C++ Contest Reference)}}
\author{Orange Juice Reference Document}
\date{\today}
\maketitle

\section{Basic}

\subsection{1.1 C++ Solution Template}
Standard competitive programming setup including Fast I/O, common header includes, DEBUG macros for variable tracing, and compilation flags for GCC.
\begin{lstlisting}[language=C++]
#include <bits/stdc++.h>
#include <stdio.h>
#include <iostream>
#include <sstream>
#include <iomanip>
#include <cstring>
#include <cmath>
#include <algorithm>
#include <climits>
#include <vector>
#include <stack>
#include <queue>
#include <set>
#include <unordered_set>
#include <map>
#include <unordered_map>
#include <cassert>

#ifdef DEBUG
#define SHOW(x) {cout << __LINE__ << ":\t"; {x} cout << endl;}
#define $(x) {cout << #x << " = " << x << ", ";}
#define REACH_HERE { printf("%d:\tREACH_HERE\n", __LINE__); }
#define PRINT(s, ...) { printf("%d:\t", __LINE__); printf(s, ##__VA_ARGS__); } 
#define PRINTLN(s, ...) { printf("%d:\t", __LINE__); printf(s, ##__VA_ARGS__); printf("\n"); }
#else
#define SHOW(x) {}
#define $(x) {}
#define REACH_HERE {}
#define PRINT(s, ...) {} 
#define PRINTLN(s, ...) {}
#endif

const double E = 1e-8;
const double PI = acos(-1);

using namespace std;

int main() {
    ios::sync_with_stdio(false);

    int a = 10, b = 10;
    SHOW( $(a) $(b) )

    PRINTLN("Hello %d", a)
}
\end{lstlisting}

\subsection{1.2 Strings}
Basic C-string and C++ \texttt{std::string} conversions and input processing tricks.
\begin{lstlisting}[language=C++]
string a;
getline(cin, a);
cout << a << endl;
char target[1024];
strcpy(target, cppstr.c_str());
char s[12];
gets(s);   // till new line or EOF.
string target = string(arr);
\end{lstlisting}

\subsection{1.3 STL Algorithm}

\subsubsection{1.3.1 Permutation}
Generates next lexicographical permutation. Useful for brute-force over small sets ($N \le 10$).
\begin{lstlisting}[language=C++]
bool next_permutation (BidirectionalIterator first, BidirectionalIterator last);
bool next_permutation (BidirectionalIterator first, BidirectionalIterator last, Compare comp);
do {
    std::cout << myints[0] << ' ' << myints[1] << ' ' << myints[2] << '\n';
} while ( std::next_permutation(myints,myints+3) );
\end{lstlisting}

\subsubsection{1.3.2 Binary Search \& Lower Bound}
\texttt{lower\_bound} returns an iterator to the first element not less than val.
\begin{lstlisting}[language=C++]
bool binary_search (ForwardIterator first, ForwardIterator last, const T& val, Compare comp);
ForwardIterator lower_bound (ForwardIterator first, ForwardIterator last, const T& val, Compare comp);
\end{lstlisting}

\subsubsection{1.3.4 Swap}
\begin{lstlisting}[language=C++]
void swap (T& a, T& b);
void iter_swap (ForwardIterator1 a, ForwardIterator2 b);
\end{lstlisting}

\subsubsection{1.3.5 Heap}
STL inline heap transformations on raw random access iterators (vectors/arrays).
\begin{lstlisting}[language=C++]
void make_heap (RandomAccessIterator first, RandomAccessIterator last, Compare comp);
void pop_heap (RandomAccessIterator first, RandomAccessIterator last, Compare comp);
void push_heap (RandomAccessIterator first, RandomAccessIterator last, Compare comp);
void sort_heap (RandomAccessIterator first, RandomAccessIterator last, Compare comp);
\end{lstlisting}

\subsubsection{1.3.6 Sort}
\begin{lstlisting}[language=C++]
void sort (RandomAccessIterator first, RandomAccessIterator last, Compare comp);
void stable_sort ( RandomAccessIterator first, RandomAccessIterator last, Compare comp );
\end{lstlisting}

\subsubsection{1.3.7 Compare Custom Implementations}
Multiple ways to pass custom comparison criteria to containers and algorithms.
\begin{lstlisting}[language=C++]
// Lambda expression
auto cmp = [](const T& a, const T& b) { return true; };
set<T, decltype(cmp)> a_set_with_customized_comparator(cmp);

// Compare function
bool cmp (int i,int j) { return (i<j); }
set<int, bool (*)(int, int)> s(cmp);

// Define operator <()
struct Edge {
   int from, to, weight;
   bool operator<(Edge that) const {
        return weight > that.weight;
    }
};

// Define operator()()
vector<int> occurrences;
struct cmp {
    bool operator()(int a, int b) {
        return occurrences[a] < occurrences[b];
    }
};
set<int, cmp> s;
priority_queue<int, vector<int>, cmp> pq;
\end{lstlisting}

\subsection{1.4 STL Containers}

\subsubsection{1.4.1 Map \& Unordered Map}
Ordered Red-Black Tree vs Hash Map usage.
\begin{lstlisting}[language=C++]
template < class Key, class T, class Compare = less<Key>, class Alloc = allocator<pair<const Key,T> > > class map;
template < class Key, class T, class Hash = hash<Key>, class Pred = equal_to<Key>, class Alloc = allocator< pair<const Key,T> > > class unordered_map;
\end{lstlisting}

\subsubsection{1.4.3 - 1.4.7 Vector, Queue, Deque, Stack}
\begin{lstlisting}[language=C++]
// Vector
std::vector<int> second (4,100);

// Queue
queue<int> my_queue;
queue<int, list<int> > my_queue (my_list);

// Stack
stack<int, vector<int> > my_stack (my_data);
\end{lstlisting}

\subsubsection{1.4.8 Priority Queue \& Array Min-Max Heap}
Priority queue initialization and custom min-max heap array implementation.
\begin{lstlisting}[language=C++]
priority_queue<int> my_priority_queue;
priority_queue<int, vector<int>, greater<int> > two_priority_queue;

struct HEAP {
     int arr[127];
     int size;
};
void insert(HEAP *h,int n){
       int sz = h->size++;
       while(sz>0){
            if(n < h->arr[(sz-1)/2]){
                h->arr[sz] = h->arr[(sz-1)/2];
                sz = (sz-1)/2;
            }else{
                h->arr[sz] = n;
                return;
            }
       }
       h->arr[0] = n;
}
void delete_heap(HEAP *h){
       h->size--;
       h->arr[0] = h->arr[h->size];
       h->arr[h->size] = 0;
       int i = 1,cv = h->arr[0];
       while(i < h->size){
               if(i < h->size-1 && h->arr[i] > h->arr[i+1])
                      i++;
               if(h->arr[i] < cv){
                      h->arr[(i-1)/2] = h->arr[i];
                      h->arr[i] = cv;
                      i = i*2 + 1;
               }else{
                      return;
               }
       }
}
\end{lstlisting}

\section{2. Advanced Data Structures}

\subsection{2.2 Tree}

\subsubsection{2.2.1 Pointer Jumping (Binary Lifting)}
Preprocessing: $O(N \log N)$, Query: $O(\log N)$. Used for finding $K$-th ancestor in a tree structure.
\begin{lstlisting}[language=C++]
#define MAX_NODE 100030
#define MAX_NODE_LOG 20
#define TREE_ROOT 0

vector<int> g[MAX_NODE];
vector<int> parent_jump[MAX_NODE];
vector<int> path;

void init_jump(int cur = TREE_ROOT) {
	int d = 1;
	while (true) {
		int index = path.size() - d;
		if (index < 0) break;
		parent_jump[cur].push_back(path[index]);
		d <<= 1;
	}
	path.push_back(cur);

    for (int i = 0; i < g[cur].size(); i++) {
        int nx = g[cur][i];
        if (cur == TREE_ROOT || nx != parent_jump[cur][0]) {
            init_jump(nx);
        }
    }
    path.pop_back();
}

int go_up(int cur, int dis) {
	int mask = 1, index = 0;
	while (mask <= dis) {
		if (dis & mask)
			cur = parent_jump[cur][index];
		mask <<= 1;
		index++;
	}
	return cur;
}
\end{lstlisting}

\subsubsection{2.2.2 Heavy-Light Decomposition (HLD)}
Decomposes a tree into paths (chains) to query path operations using a segment tree in $O(\log^2 N)$.
\begin{lstlisting}[language=C++]
struct Graph {
    map<pair<int, int>, int> node_edge;
    struct Edge { int from, to; unsigned long long y; };
    const static int MAXNODE = 200005;
    vector<int> g[MAXNODE];
    vector<Edge> edge;
    int n, parent[MAXNODE], size[MAXNODE];

    void init(int nn) {
        n = nn;
        for (int i = 0; i <= n; i++) g[i].clear();
        edge.clear();
    }

    void add_e(int x, int y, unsigned long long val_y) {
    	Edge e1 = {x, y, val_y};
        g[x].push_back(edge.size());
        node_edge[make_pair(x, y)] = edge.size();
        edge.push_back(e1);
        Edge e2 = {y, x, val_y};
        g[y].push_back(edge.size());
        node_edge[make_pair(y, x)] = edge.size();
        edge.push_back(e2);
    }

    int size_parent(int cur) {
        int& s = size[cur] = 1;
        for (int i = 0; i < g[cur].size(); i++) {
            int ie = g[cur][i];
            const Edge& e = edge[ie];
            int nx = e.to;
            if (nx != parent[cur]) {
                parent[nx] = cur;
                s += size_parent(nx);    
            }
        }
        return s;
    }

    struct Chain { int head, head_depth, len; };
    vector<Chain> chain;
    int chain_total, chain_no[MAXNODE], chain_pos[MAXNODE];

    void hld(int cur, int depth) {
        if (chain_total == chain.size()) {
            Chain c = {cur, depth, 0};
            chain.push_back(c);
        }
        chain_no[cur] = chain_total;
        chain_pos[cur] = chain[chain_total].len++;
        if (depth != 0 && g[cur].size() - 1 == 0) return;

        int heavy = edge[g[cur][0]].to;
        for (int i = 0; i < g[cur].size(); i++) {
            int nx = edge[g[cur][i]].to;
            if (heavy == parent[cur] || (nx != parent[cur] && size[nx] > size[heavy]))
                heavy = nx;
        }

        hld(heavy, depth + 1);
        for (int i = 0; i < g[cur].size(); i++) {
            int nx = edge[g[cur][i]].to;
            if (nx != parent[cur] && nx != heavy) {
                chain_total++;
                hld(nx, depth + 1);
            }
        }
    }

    void heavy_light_decomposition() {
        int root = 1; parent[root] = -1;
        size_parent(root);
        chain_total = 0; chain.clear();
        hld(root, 0); chain_total++;
    }
};
\end{lstlisting}

\subsubsection{2.2.3 Lowest Common Ancestor (LCA using RMQ)}
Euler tour + RMQ technique to compute LCA in $O(1)$ after $O(N \log N)$ preprocessing.
\begin{lstlisting}[language=C++]
struct Graph {
    struct Edge { int to, len; };
    const static int MAXNODE = 100002;
    vector<int> g[MAXNODE];
    vector<Edge> edge;
    int n, root = 0;

    vector<int> dis_to_root, first_visit_time, visit;
    int visit_counter;
    vector<vector<int>> rmq;

    int range_minimum_query(int l, int r) {
        if (l > r) swap(l, r);
        int interval_len = r - l;
        int first_half = 1;
        while ((1 << first_half) <= interval_len) first_half++;
        first_half--;
        int second_half = r - (1 << first_half) + 1;
        if (first_visit_time[rmq[l][first_half]] < first_visit_time[rmq[second_half][first_half]])
            return rmq[l][first_half];
        return rmq[second_half][first_half];
    }

    int get_lca(int x, int y) {
        return range_minimum_query(first_visit_time[x], first_visit_time[y]);
    }

    void euler_tour(int cur) {
        visit[++visit_counter] = cur;
        if (first_visit_time[cur] == 0) first_visit_time[cur] = visit_counter;
        for (int ie : g[cur]) {
            int nx = edge[ie].to, len = edge[ie].len;
            if (first_visit_time[nx] == 0) {
                dis_to_root[nx] = dis_to_root[cur] + len;
                euler_tour(nx);
                visit[++visit_counter] = cur;
            }
        }
    }
};
\end{lstlisting}

\subsubsection{2.2.3.1 Tarjan's Off-line Algorithm for LCA}
Offline algorithm to process all LCA queries in $O(N + M)$ using Disjoint Set Union (DSU).
\begin{lstlisting}[language=C++]
#define MAXHHH 50003
#define MAXJJJ 75005

struct Node {
    vector<int> next, dist, query, lca, q_i;
};
Node g[MAXHHH];
int father[MAXHHH], came[MAXHHH];

int find(int x) { 
    return father[x] == x ? x : father[x] = find(father[x]);
}
void mergeFirstInToSecond(int x, int y) { father[find(x)] = find(y); }

void tarjan_lca_dfs(int cur) {
    came[cur] = 1;
    for (unsigned int i = 0; i < g[cur].next.size(); i++) {
        int next = g[cur].next[i];
        if (came[next] == 1) continue;
        tarjan_lca_dfs(next);
        mergeFirstInToSecond(cur, next);
    }
    for (unsigned int i = 0; i < g[cur].query.size(); i++) {
        int that = g[cur].query[i];
        if (came[that] == 0) continue;
        g[cur].lca[i] = find(that);
    }
}
\end{lstlisting}

\subsubsection{2.2.4 Centroid Decomposition}
Divide and conquer technique on trees. Useful for path queries passing through tree centroids.
\begin{lstlisting}[language=C++]
int dfs_size(int cur, int from) {
    int total = 1;
    for (int child : tree[cur].children) {
        if (child == from) continue;
        total += dfs_size(child, cur);
    }
    tree[cur].size = total;
    return total;
}

void find_centroid(int cur, int from) {
    if (tree[cur].size == 1) tree[cur].c_parent = from;
    int max_child_size = 0, max_child;
    for (int child : tree[cur].children) {
        if (tree[child].c_parent != -1) continue;
        if (tree[child].size > max_child_size) {
            max_child_size = tree[child].size;
            max_child = child;
        }
    }
    if (max_child_size > tree[cur].size / 2) {
        tree[max_child].size = tree[cur].size;
        tree[cur].size -= max_child_size;
        find_centroid(max_child, from);
    } else {
        tree[cur].c_parent = from;
        if (from == 0) ctree_root = cur;
        for (int child : tree[cur].children) {
            if (tree[child].c_parent != -1) continue;
            find_centroid(child, cur);
        }
    }
}
\end{lstlisting}

\subsection{2.3 Trie / Trie Graph / AC Automaton}
Aho-Corasick Automaton for multi-pattern string matching in $O(N_{len} + M_{article})$.
\begin{lstlisting}[language=C++]
#define HHH 1000002
struct TrieNode {
    char val; bool ended; int count, childCount;
    int next[26], prev; bool suffixEnded; int suffix[26];
    TrieNode() {
        val = 0; memset(next, -1, sizeof(next));
        ended = false; count = childCount = 0;
        prev = -1; suffixEnded = false; memset(suffix, -1, sizeof(suffix));
    }
};

struct Trie {
    TrieNode node[HHH]; int size;
    int add(string& s) {
        int preIndex = 0;
        for (int i = 0; i < s.length(); i++) {
            int& curIndex = node[preIndex].next[s[i] - 'a'];
            if (curIndex == -1) curIndex = ++size;
            node[curIndex].val = s[i];
            node[curIndex].childCount++;
            preIndex = curIndex;
        }
        node[preIndex].ended = true; node[preIndex].count++;
        return preIndex;
    }

    void buildSuffix() {
        queue<int> q;
        for (int i = 0; i < 26; i++) {
            int& next = node[0].next[i];
            node[0].suffix[i] = 0;
            if (next == -1) next = 0;
            else { q.push(next); node[next].prev = 0; }
        }
        while (!q.empty()) {
            int cur = q.front(); q.pop();
            int prevSuffix = node[node[cur].prev].suffix[node[cur].val - 'a'];
            if (node[prevSuffix].ended) node[cur].suffixEnded = true;
            for (int i = 0; i < 26; i++) {
                int& next = node[cur].next[i];
                node[cur].suffix[i] = node[prevSuffix].next[i];
                if (next == -1) next = node[cur].suffix[i];
                else { q.push(next); node[next].prev = cur; }
            }
        }
    }
};
\end{lstlisting}

\subsection{2.5 Suffix Array}
Suffix Array construction in $O(N \log N)$ and LCP computation in $O(N)$.
\begin{lstlisting}[language=C++]
#define HH 100002
string s;
int rank_array[HH], rank_array_temp[HH], suffix_array[HH], suffix_array_temp[HH], counter[HH];

void counting_sort(int k) {
	memset(counter, 0, sizeof(counter));
	int len = s.length();
	for (int i = 0; i < len; i++) {
		int old_rank = i + k < len ? rank_array[i + k] : 0;
		counter[old_rank]++;
	}
	int accu = 0, largest_possible_value = max(256, len);
	for (int i = 0; i < largest_possible_value; i++) {
		int count_temp = counter[i];
		counter[i] = accu;
		accu += count_temp;
	}
	for (int i = 0; i < len; i++) {
		int old_rank = suffix_array[i] + k < len ? rank_array[suffix_array[i] + k] : 0;
		suffix_array_temp[counter[old_rank]++] = suffix_array[i];
	}
	memcpy(suffix_array, suffix_array_temp, sizeof(suffix_array));
}

void build_suffix_array() {
	int len = s.length();
	for (int i = 0; i < len; i++) rank_array[i] = s[i];
	for (int i = 0; i < len; i++) suffix_array[i] = i;
	for (int k = 1; k < len; k <<= 1) {
		counting_sort(k); counting_sort(0);
		int rank = 0;
		rank_array_temp[suffix_array[0]] = rank;
		for (int i = 1; i < len; i++) {
			if (rank_array[suffix_array[i - 1]] != rank_array[suffix_array[i]] ||
				rank_array[suffix_array[i] + k] != rank_array[suffix_array[i - 1] + k])
				rank++;
			rank_array_temp[suffix_array[i]] = rank;
		}
		memcpy(rank_array, rank_array_temp, sizeof(rank_array));
	}
}
\end{lstlisting}

\subsubsection{2.5.3 Longest Common Prefix (LCP)}
\begin{lstlisting}[language=C++]
int longest_common_prefix[HH], phi[HH], permuted_lcp[HH];

void compute_lcp() {
	int len = s.length();
	phi[suffix_array[0]] = -1;
	for (int i = 1; i < len; i++) phi[suffix_array[i]] = suffix_array[i - 1];
	for (int i = 0, cur_lcp = 0; i < len; i++) {
		if (phi[i] == -1) permuted_lcp[i] = 0;
		else {
			while (s[i + cur_lcp] == s[phi[i] + cur_lcp]) cur_lcp++;
			permuted_lcp[i] = cur_lcp;
			cur_lcp = max(cur_lcp - 1, 0);
		}
	}
	for (int i = 0; i < len; i++) longest_common_prefix[i] = permuted_lcp[suffix_array[i]];
}
\end{lstlisting}

\subsection{2.6 Binary Indexed Tree (Fenwick Tree)}
Point updates and range prefix sum queries in $O(\log N)$.
\begin{lstlisting}[language=C++]
#define MAX_INDEX 100005
int tree[MAX_INDEX + 1];

int low_bit(int i) { return i & -i; }

int query(int i) {
    int ans = 0;
    for (; i > 0; i -= low_bit(i)) ans += tree[i];
    return ans;
}

void insert(int i, int value) {
    for (; i <= MAX_INDEX; i += low_bit(i)) tree[i] += value;
}
\end{lstlisting}

\subsection{2.7 Segment Tree (Lazy Propagation)}
Range Query and Range Update in $O(\log N)$.
\begin{lstlisting}[language=C++]
struct SegmentTree {
    struct Op { int h; };
    struct Node { int l, r, h; bool lazy; Op op; };
    vector<Node> node;

    void init(int l, int r) {
        int tree_range = r - l + 1, tree_size = 1;
        while (tree_size <= tree_range) tree_size <<= 1;
        tree_size <<= 1;
        node.resize(tree_size);
        node[1] = {l, r, 0, false, {0}};
    }

    int query(int xl, int xr, int i = 1) {
        Node& cur = node[i];
        if (cur.l == cur.r || (xl <= cur.l && cur.r <= xr)) return cur.h;
        int lci = i * 2, rci = lci + 1;
        if (cur.lazy) {
            update(node[lci].l, node[lci].r, cur.op.h, lci);
            update(node[rci].l, node[rci].r, cur.op.h, rci);
            cur.lazy = false;
        }
        int ret = INT_MAX;
        if (xl <= node[lci].r) ret = min(ret, query(xl, xr, lci));
        if (node[rci].l <= xr) ret = min(ret, query(xl, xr, rci));
        return ret;
    }

    void update(int xl, int xr, int xh, int i = 1) {
        Node& cur = node[i];
        if (cur.l == cur.r || (xl <= cur.l && cur.r <= xr)) {
            cur.h = max(cur.h, xh);
            cur.op.h = max(cur.op.h, xh);
            cur.lazy = true;
            return;
        }
        int lci = i * 2, rci = lci + 1;
        if (xl <= node[lci].r) update(xl, xr, xh, lci);
        if (node[rci].l <= xr) update(xl, xr, xh, rci);
        cur.h = min(node[lci].h, node[rci].h);
    }
};
\end{lstlisting}

\subsection{2.8 Range Minimum Query RMQ (Sparse Table)}
Static Range Minimum Queries in $O(1)$ query time after $O(N \log N)$ preprocessing.
\begin{lstlisting}[language=C++]
struct RMQ {
    const static int MAXLENGTH = 200005, LOG_MAXLENGTH = 20;
    int rmq[MAXLENGTH][LOG_MAXLENGTH];

    void init(int* arr, int len) {
        for (int i = 0; i < len; i++) rmq[i][0] = arr[i];
        for (int j = 1; j < LOG_MAXLENGTH; j++)
            for (int i = 0; i + (1 << j) <= len; i++)
                rmq[i][j] = min(rmq[i][j - 1], rmq[i + (1 << (j - 1))][j - 1]);
    }

    int range_minimum_query(int l, int r) {
        if (l > r) swap(l, r);
        int k = 32 - __builtin_clz(r - l + 1) - 1;
        return min(rmq[l][k], rmq[r - (1 << k) + 1][k]);
    }
};
\end{lstlisting}

\subsection{2.9 Union-find Set (Disjoint Set Union)}
\begin{lstlisting}[language=C++]
struct UnionFindSet {
	vector<int> parent;
	void init(int nn) {
		parent.resize(nn + 1);
		iota(parent.begin(), parent.end(), 0);
	}
	void merge(int x, int y) { parent[find(x)] = find(y); }
	int find(int x) { return x == parent[x] ? x : parent[x] = find(parent[x]); }
	bool together(int x, int y) { return find(x) == find(y); }
};
\end{lstlisting}

\section{3. Methodology}

\subsection{3.1 Recursive: Hanoi}
\begin{lstlisting}[language=C++]
void hanoi(int n, char x, char y, char z) {
    if (n == 1) printf("%d from %c to %c\n", n, x, z);
    else {
        hanoi(n-1, x, z, y);
        printf("%d from %c to %c\n", n, x, z);
        hanoi(n-1, y, x, z);
    }
}
\end{lstlisting}

\subsection{3.2 Dynamic Programming: LIS}
Longest Increasing Subsequence in $O(N \log N)$ using binary search.
\begin{lstlisting}[language=C++]
vector<int> sequence;
vector<int> lis(sequence.size() + 1, INT_MAX);
for (int i = 0; i < sequence.size(); i++) {
    int r = sequence[i];
    auto ptr = lower_bound(begin(lis), end(lis), r);
    *ptr = min(*ptr, r);
}
\end{lstlisting}

\section{4. String}

\subsection{4.1 KMP Algorithm}
String pattern matching in $O(N + M)$ time complexity.
\begin{lstlisting}[language=C++]
#define HHH 10003
int ne[HHH];

int kmp(string& par, string& ori) {
    ne[0] = -1;
    for (int p = ne[0], i = 1; i < par.length(); i++) {
        while (p >= 0 && par[p+1] != par[i]) p = ne[p];
        if (par[p+1] == par[i]) p++;
        ne[i] = p;
    }
    int match = 0;
    for (int p = -1, q = 0; q < ori.length(); q++) {
        while (p >= 0 && par[p+1] != ori[q]) p = ne[p];
        if (par[p+1] == ori[q]) p++;
        if (p + 1 == par.length()) { p = ne[p]; match++; }
    }
    return match;
}
\end{lstlisting}

\subsection{4.3 Manacher's Algorithm}
Finds the longest palindromic substring in $O(N)$ linear time.
\begin{lstlisting}[language=C++]
int dp[10003];
int lengthLongestPalindromSubstring(string& s) {
    memset(dp, 0, sizeof(dp));
    int ans = 0, pivot = 1, len = s.length() * 2;
    for (int i = 1; i < len; i++) {
        int pBorder = pivot + dp[pivot], iBorder = i;
        if (iBorder < pBorder && 2 * pivot - i > 0) {
            dp[i] = dp[2*pivot-i];
            iBorder = min(pBorder, i + dp[i]);
        }
        if (iBorder >= pBorder) {
            int j = iBorder + (iBorder % 2 ? 2 : 1);
            for (; j < len && 2*i-j > 0 && s[j/2] == s[(2*i-j)/2]; j += 2);
            iBorder = j - 2; dp[i] = iBorder - i; pivot = i;
        }
        ans = max(ans, dp[i] + 1);
    }
    return ans;
}
\end{lstlisting}

\section{5. Graph}

\subsection{5.2 Minimum Spanning Tree}

\subsubsection{5.2.1 Prim's Algorithm}
$O((V + E) \log V)$ for finding MST using Priority Queue.
\begin{lstlisting}[language=C++]
vector<bool> added; vector<int> mindis;
void mst() {
    vector<bool>(n + 1, false).swap(added);
    vector<int>(n + 1, INT_MAX).swap(mindis);
    auto cmp = [](const Edge& a, const Edge& b) { return a.len > b.len; };
    priority_queue<Edge, vector<Edge>, decltype(cmp)> near(cmp);
    for (int ie : g[1]) { near.push(edge[ie]); mindis[edge[ie].to] = edge[ie].len; }
    added[1] = true;
    while (!near.empty()) {
        Edge cur = near.top(); near.pop();
        if(added[cur.to]) continue;
        added[cur.to] = true;
        for (int ie : g[cur.to]) {
            int nx = edge[ie].to;
            if (!added[nx] && mindis[nx] > edge[ie].len) {
                mindis[nx] = edge[ie].len;
                near.push(edge[ie]);
            }
        }
    }
}
\end{lstlisting}

\subsubsection{5.2.2 Kruskal's Algorithm}
$O(E \log E)$ MST using Disjoint-Set Union.
\begin{lstlisting}[language=C++]
UnionFindSet ufs;
void mst() {
    ufs.init(n);
    vector<Edge> eee = edge;
    sort(begin(eee), end(eee), [](const Edge& a, const Edge& b) { return a.len < b.len; });
    int need = n - 1;
    for (const auto& e : eee) {
        if (!ufs.together(e.from, e.to)) {
            ufs.merge(e.from, e.to);
            if (!(--need)) break;
        }
    }
}
\end{lstlisting}

\subsection{5.3 Shortest Path}

\subsubsection{5.3.4 Dijkstra's Algorithm}
Shortest path algorithm on non-negative weighted graphs in $O(E \log V)$.
\begin{lstlisting}[language=C++]
struct cmp {
    bool operator()(pair<int, int> a, pair<int, int> b) { return a.first > b.first; }
};
void dijkstra(int s) {
    priority_queue<pair<int, int>, vector<pair<int, int>>, cmp> q;
    q.push(make_pair(0, s)); dist[s] = 0;
    while (!q.empty()) {
        int cur = q.top().second; q.pop();
        if (done[cur]) continue;
        done[cur] = true;
        for (int next : children[cur]) {
            if (!done[next] && dist[next] > dist[cur] + edge[cur][next]) {
                dist[next] = dist[cur] + edge[cur][next];
                q.push(make_pair(dist[next], next));
            }
        }
    }
}
\end{lstlisting}

\subsection{5.4 Maximum Matching}

\subsubsection{5.4.1.1 Hungarian Algorithm}
Finds maximum cardinality matching in bipartite graphs in $O(V \cdot E)$.
\begin{lstlisting}[language=C++]
int peer[405]; bool went[405];
bool match(int cur) {
    for (int ie = last[cur]; ie != -1; ) {
        const Edge& e = edge[ie]; ie = e.pre_edge; int to = e.to;
        if (went[to]) continue;
        went[to] = true;
        if (peer[to] == -1 || match(peer[to])) {
            peer[to] = cur; peer[cur] = to;
            return true;
        }
    }
    return false;
}
int bipartite_match() {
    fill(peer, peer + nv, -1); int ans = 0;
    for (int i = 0; i < nv; i++) {
        if (last[i] == -1 || peer[i] != -1) continue;
        fill(went, went + nv, false);
        if (match(i)) ans++;
    }
    return ans;
}
\end{lstlisting}

\subsubsection{5.4.2.1 Blossom Algorithm}
Maximum matching on general non-bipartite graphs.
\begin{lstlisting}[language=C++]
const int NMax = 230;
int Next[NMax], spouse[NMax], belong[NMax], Q[NMax], bot, mark[NMax], visited[NMax];
vector<int> E[NMax]; int N;

int findb(int a) { return belong[a] == a ? a : belong[a] = findb(belong[a]); }
void together(int a,int b) {
    a = findb(a); b = findb(b);
    if (a != b) belong[a] = b;
}

int findLCA(int x,int y) {
    static int t = 0; t++;
    while (1) {
        if (x != -1) {
            x = findb(x);
            if (visited[x] == t) return x;
            visited[x] = t;
            x = (spouse[x] != -1) ? Next[spouse[x]] : -1;
        }
        swap(x,y);
    }
}

void goup(int a,int p) {
    while (a != p) {
        int b = spouse[a], c = Next[b];
        if (findb(c) != p) Next[c] = b;
        if (mark[b] == 2) mark[Q[bot++] = b] = 1;
        if (mark[c] == 2) mark[Q[bot++] = c] = 1;
        together(a,b); together(b,c); a = c;
    }
}

void findaugment(int s) {
    for (int i = 0; i < N; i++) Next[i] = -1, belong[i] = i, mark[i] = 0, visited[i] = -1;
    Q[0] = s; bot = 1; mark[s] = 1;
    for (int head = 0; spouse[s] == -1 && head < bot; head++) {
        int x = Q[head];
        for (int y : E[x]) {
            if (spouse[x] != y && findb(x) != findb(y) && mark[y] != 2) {
                if (mark[y] == 1) {
                    int p = findLCA(x,y);
                    if (findb(x) != p) Next[x] = y;
                    if (findb(y) != p) Next[y] = x;
                    goup(x,p); goup(y,p);
                } else if (spouse[y] == -1) {
                    Next[y] = x;
                    for (int j = y; j != -1; ) {
                        int k = Next[j], l = spouse[k];
                        spouse[j] = k; spouse[k] = j; j = l;
                    }
                    break;
                } else {
                    Next[y] = x;
                    mark[Q[bot++] = spouse[y]] = 1; mark[y] = 2;
                }
            }
        }
    }
}
\end{lstlisting}

\end{document}